\section{另类架构、1T 之后与模型进化史}

访谈后半段回顾过去三年的模型进化：ChatGPT 释放短上下文聊天交互；2023 年开源团队追赶预训练；reasoning 阶段代码和数学成为可验证信号；Agent 阶段代码和环境反馈成为长程任务训练入口。罗福莉认为，国内团队在 pre-train 上代差已经很小，下一步竞争在 Agent post-train 和 RL scaling。

\begin{figure}[H]
\centering
\includegraphics[width=0.96\textwidth]{model-evolution-timeline.png}
\caption{过去三年模型进化：从 ChatGPT 到 Agent 后训练。}
\end{figure}

\begin{knowledgebox}{读图：每一阶段如何继承上一阶段}
时间线不是替代关系。ChatGPT 让用户感知预训练智能；开源追赶补齐数据和架构；reasoning 让可验证任务成为训练信号；Agent 原型把代码和环境引入长程任务；2026 年则把 post-train scaling 推到主线。
\end{knowledgebox}

\subsection{另类架构的意义}

罗福莉提到 DeepSeek-V2、MoE、attention 改造等“反主流”选择。另类架构不是为了反主流而反主流，而是在资源受限或成本敏感环境中寻找可扩展研究。一个研究若只是 paper trick，不一定能转化为工业级模型；但若能 scale，并最终形成高水平模型状态，就是有价值的 frontier research。

\noindent\begin{tabularx}{\textwidth}{p{0.22\textwidth}p{0.34\textwidth}X}
\toprule
路线 & 解决的问题 & 工业化判断 \\
\midrule
MoE & 用稀疏激活降低单位 token 计算成本 & 需要路由、负载均衡、训练稳定性和部署效率同时成立。 \\
Attention 改造 & 降低长上下文和推理成本 & 要看能否在真实任务中保持质量，而非只看局部 benchmark。 \\
小模型 + 框架 & 用便宜模型配合 Agent scaffold & 适合高频、低成本、可被框架补短板的环节。 \\
1T 基座 & 提供前沿 Agent 能力底座 & 是入场券，不是完整答案。 \\
\bottomrule
\end{tabularx}

\subsection{AI 没有生存危机}

罗福莉用人类进化和模型进化做对比：人类是在自然环境和生存压力中进化，模型没有自己的生存危机，却拥有大量算力、人类知识和人类帮助。因此 AI 的进化路径可能更自由、更散漫、更快，也更不同于生物智能。这一段提醒我们，不要把人类进化类比生硬套到模型上。

\begin{knowledgebox}{类比的边界}
“语言在人类智能中很晚出现”可以启发我们理解 AI 的倒三角结构，但模型环境与人类环境完全不同。AI 的下一步不一定复刻人类身体或感官路径，而可能先在 coding、生产力和数字环境中继续扩张。
\end{knowledgebox}

\subsection{每天否认昨天的自己}

罗福莉说，过去半年几乎每天都在否认昨天的自己。这个表述很适合 2026 年 Agent 阶段：工具、框架、模型、成本和组织方式都在快速变化。对研究团队来说，稳定不是固守旧判断，而是建立能快速更新判断的机制。

\begin{importantbox}{研究速度的来源}
研究速度不只来自聪明人，也来自环境：足够算力、足够好的基础设施、足够快的反馈、足够开放的组织和足够明确的价值观。环境比经验更重要，是这期结尾的关键判断。
\end{importantbox}

\subsection{本章小结}

过去三年的模型进化，从聊天交互、开源追赶、reasoning 验证，走到 Agent 后训练。下一阶段的领先取决于可 scale 的架构、后训练系统、资源配置和快速自我否定能力。

